

%\DeclareFontFamily{OT1}{pzc}{}
%\DeclareFontShape{OT1}{pzc}{m}{it}{<->s*[1.20]pzcmi7t}{}
\DeclareMathAlphabet{\mathpzc}{OT1}{pzc}{m}{it}


\mathcode`\:="603A     % : = \colon
\mathchardef\colon="303A  % :=
\newcommand{\typing}{\colon}

\newcommand{\define}[1]{\emph{#1}}
\newcommand{\arcmp}{\circ}
\newcommand{\celcmp}{*}
\newcommand{\rwhisk}[1]{#1} %writes the 1-cell to the right
\newcommand{\lwhisk}[1]{#1} %writes the 2-cell to the right
\newcommand{\mate}[1]{#1^{\#}} 

\def\ct#1{\ensuremath{\mathpzc{#1}}}
\newcommand{\fun}[1]{\mathrm{#1}}
\def\cttwo#1{\ensuremath{\mathbf{#1}}}
\def\op{\ensuremath{^{\mathrm{op}}}}
\def\dfn#1{{\it #1}}
\def\rstn{\ensuremath{\mathord{\restriction}}}
\def\gint{\ensuremath{\mathord{\int}\kern-.2ex}}
\def\Hom#1{\ensuremath{#1}}
\def\OPR#1{\ensuremath{\mathop{\mathrm{#1}}\nolimits}} % perché \mathop?
\def\Obj{\OPR{Obj}}
\def\Mor{\OPR{Mor}}
\def\Open{\OPR{Open}}
\def\Aut{\ct{Aut}}
\def\Arr#1{\ensuremath{#1^{2}}}
\def\cod{\OPR{cod}}
\def\dom{\OPR{dom}}
\def\Gint#1{\gint\!\gint #1}
\def\id#1{\ensuremath{\mathrm{id}_{#1}}} %identity 1-cells
\def\idtwo#1{\ensuremath{\mathrm{i}_{#1}}} %identity 2-cells
\def\Id#1{\ensuremath{\mathrm{Id}_{#1}}} %identity functors between categories
\def\twoId#1{\ensuremath{\boldsymbol{\mathrm{Id}_{#1}}}} %identity functors between 2-categories
\newcommand{\zero}{\ensuremath{0}}
\newcommand{\obzero}[1]{\ensuremath{0}^#1}
\newcommand{\coprodu}{+}

\def\fp#1{\ensuremath{|#1|}}  %components of arrows in jacobs completion
\def\sp#1{\ensuremath{\num{#1}}}
\def\tp#1{\ensuremath{\widetilde{#1}}}

\def\ini#1#2{\ensuremath{#1\hookrightarrow #2}}  %initial segment inclusion
\def\carr#1{\ensuremath{#1^+}}  %comprehension arrow
\def\Cat{\ensuremath{\ct{Cat}}}  %categoria cat
\def\Cattwo{\ensuremath{\cttwo{Cat}}}  %2-categoria cat
\def\SCat{\ensuremath{\ct{SCat}}}  %categoria semicat

\newcommand{\ctset}{\ct{Set}}  %categoria set
\newcommand{\ctsetp}{\ct{Set}_*}  %categoria set puntata
\newcommand{\Mod}[1]{#1^{*}}  %funtore mod su frecce
\newcommand{\simp}[1]{\mathrm{s(\ct{#1})}}  %simple fibration
\newcommand{\simpfib}[1]{\mathrm{s_{\ct{#1}}}}
\newcommand{\Fam}[1]{\mathrm{Fam}(#1)}
\newcommand{\Famfib}[1]{\mathrm{Fam}_{#1}}
\newcommand{\pw}[1]{\mathcal{P}(#1)}  %powerset
\newcommand{\pwn}{\pw{\mathbb{N}}}  %powerset numeri naturali
\newcommand{\pwfun}{\mathcal{P}}  %powerset functor
\newcommand{\tfib}[1]{\dot{#1}}  %term fibration
\newcommand{\tcat}[1]{\dot{#1}}  %term category
\newcommand{\tsigma}{\dot{\Sigma}}  %term sigma
\newcommand{\gtfib}[1]{\tilde{#1}}  %global type morphism fibration
\newcommand{\gtcat}[1]{\tilde{#1}}  %global type morphism category
\newcommand{\gtsigma}{\tilde{\Sigma}}  %global type morphism sigma

\newcommand{\ftransp}{\#}  %transposition functor on fibrations
\newcommand{\ttransp}{\#}  %transposition functor on total categories

\def\tdot{\textbf{.}}
%\newcommand{\nt}[3]{\ensuremath{{#1}:{#2} \stackrel{\makebox{\kern-.3ex\tdot}}\rightarrow{#3}}}
\newcommand{\nt}{\Rightarrow} %natural transformations
\newcommand{\md}{\Rrightarrow} %modifications
\newcommand{\oneAr}[3]{#1:#2\to #3} 
\newcommand{\twoAr}[3]{#1:#2\Rightarrow #3} 
\newcommand{\twofun}[1]{\mathbf{#1}}
\newcommand{\twont}[1]{\boldsymbol{#1}}
\newcommand{\prtl}{\to} %partial functions
%\newcommand{\forg}[1]{\twofun{U}_{#1}} %forgetful functor
\newcommand{\forgT}{\twofun{U_T}} %forgets terminal JTComp -> JComp
\newcommand{\forgLE}{\twofun{U_{LE}}} %forgets LE LECC -> JTComp
\newcommand{\forgJ}{\twofun{U_J}} %forgets J JComp -> Fib
\newcommand{\forgJT}{\twofun{U_{JT}}} %forgets JTComp -> Fib




\newcommand{\one}{\textbf{1}} %terminal category
\newcommand{\One}[1]{\textbf{!}_{\ct{#1}}} %terminal fibration
\newcommand{\disc}[1]{\textbf{Disc(\ct{#1})}}
\newcommand{\reind}[1]{#1^*} %reindex functor along argument
\newcommand{\carar}[2]{#1^{#2}} %cartesian lifting of the argument 1 at the argument 2

\newcommand{\num}[1]{\underline{#1}}



\def\Fib{\ensuremath{\cttwo{Fib}}}
\newcommand{\Fibone}{\ct{Fib}}
\newcommand{\Fibdisc}{\ct{Fib}_d}
\newcommand{\fpfib}{\textbf{FPFib}}
\newcommand{\fb}[1]{{#1}^{\mathrm{b}}} %freccia sotto in quadrati
\newcommand{\ft}[1]{{#1}^{\mathrm{t}}} %freccia sopra in quadrati
\newcommand{\chb}[2]{#2^\star #1} %change of base
\newcommand{\fib}[1]{#1^{-1}} 


\def\FFFP#1{\ensuremath{\ct{FFP}(#1)}}  %category

\def\fffp#1{\ensuremath{\mathbf{ffp}(#1)}}  %fibration
\def\Fffp{\ensuremath{\mathbf{ffp}}}  %functor, but actually used for the fibration
\newcommand{\Fp}{\ensuremath{\mathbf{ffp}}}  %2-functor


\def\LComp{\ensuremath{\cttwo{LECC}}} %Lawvere comprehension category
\newcommand{\Lfp}{\cttwo{FPLComp}} %LECC with fibred products

\def\JComp{\ensuremath{\cttwo{JCC}}} %Jacobs comprehension category
\def\JPComp{\ensuremath{\cttwo{JCC}_{\OPR{P}}}} %Jacobs comp. cat. w. fib. prod.
\def\JTComp{\ensuremath{\cttwo{JCC_T}}} %Jacobs comp. cat. w. fib. term.


\newcommand{\unitty}{\ensuremath{\top}}
\newcommand{\unittr}{\ensuremath{*}}
\newcommand{\type}{\ensuremath{\mathrm{Type}}}
\newcommand{\ctx}{\ensuremath{\mathrm{Ctx}}}


\def\LFib#1{\ensuremath{\hat{#1}}}  %fibrazione Lawvere libera
\def\LFun#1{\ensuremath{\hat{#1}}}  %azione del funtore libero su frecce
\def\Lun#1{\ensuremath{\twont{\eta}_{#1}}}  %funtore sopra della pseudo unità
\def\Lcoun#1{\ensuremath{\twont{\epsilon}_{#1}}}  %funtore sopra della pseudocounità
\def\LFree{\ensuremath{\twofun{\wedge}^L}}  %funtore libero
\def\LMonad{\ensuremath{\twofun{\mathrm{T}}}}  %induced monad Lawvere comprehension
\def\LTF#1{\ensuremath{\hat{#1}}}   %funtore terminale libero
\def\LCF#1{\ensuremath{\hat{#1}}}   %funtore comprensione libero
\newcommand{\term}[1]{\ensuremath{\mathrm{T}^{#1}}} %funtore terminale fibrazione
\newcommand{\comp}[1]{\ensuremath{\mathrm{C}^{#1}}} %funtore comprensione fibrazione
\newcommand{\coun}[1]{\ensuremath{\mathrm{\epsilon}^{#1}}} %counità in aggiunzione di Lawvere
\newcommand{\un}[1]{\ensuremath{\mathrm{\eta}^{#1}}} %unità in aggiunzione di Lawvere
\newcommand{\fpr}[2]{\ensuremath{\mathrm{pr}_{#1}}} %prima proiezione fibrata
\newcommand{\spr}[2]{\ensuremath{\mathrm{pr}_{#2}}}  %seconda proiezione fibrata
\newcommand{\fprj}[2]{\ensuremath{\mathrm{\pi}_{#1}}}
\newcommand{\sprj}[2]{\ensuremath{\mathrm{\pi}_{#2}}}
\newcommand{\comon}[1]{\OPR{T}}  %fibred comonad of products
\newcommand{\coKl}[1]{\ensuremath{\ct{E}}_{#1}}  %coKleisli category
\newcommand{\coKlfib}[1]{\hat{#1}}  %coKleisli fibration
\newcommand{\Kforg}[1]{\OPR{U}_{#1}}  %Kleisli forgetful functor
\def\LDom#1{\coKl{\comon{#1}}}  %dominio fibrazione Lawvere libera

\newcommand{\LtoJ}{\ensuremath{\twofun{LEJ}}}

\def\JDom#1#2{\ensuremath{{#1}^{\twofun{J}}_{#2}}}  %dominio fibrazione Jacobs libera
\def\JPDom#1#2{\ensuremath{{#1}^{#2}}}  %domain free comp. cat. w. fib. prod.
\def\JTDom#1#2{\ensuremath{{#1}^{\twofun{T}}_{#2}}}  %domain free comp. cat. w. fib. prod.
\def\JDomE#1{\JDom{\ct{E}}{#1}}  %dominio fibrazione Jacobs libera
\def\JPDomE#1{\JPDom{\ct{E}}{#1}}  %domain free comp. cat. w. fib. prod.
\def\JTDomE#1{\JTDom{\ct{E}}{#1}}  %domain free comp. cat. w. fib. prod.
\def\JFib#1{\ensuremath{{#1}^{\twofun{J}}}}  %fibrazione Jacobs libera
\def\JPFib#1{\ensuremath{\hat{#1}}}  %free comprehension preserves fib. prod.
\def\JTFib#1{\ensuremath{{#1}^{\twofun{T}}}}  %free comprehension preserves fib. term.
\def\JFun#1{\ensuremath{{#1}^{\twofun{J}}}}  %azione del funtore libero su frecce
\def\JPFun#1{\ensuremath{\hat{#1}}}
\def\JTFun#1{\ensuremath{{#1}^{\twofun{T}}}}
\def\Jun#1{\ensuremath{\twont{\eta}^{\twofun{J}}_{#1}}}  
\def\JTun#1{\ensuremath{\twont{\eta}^{\twofun{T}}_{#1}}}  
\def\Jcoun#1{\ensuremath{\twont{\epsilon}^{\twofun{J}}_{#1}}}
\def\JTcoun#1{\ensuremath{\twont{\epsilon}^{\twofun{T}}_{#1}}}
\def\JFree{\ensuremath{\twofun{F^J}}}  %funtore libero da Fib a JComp
\def\JPFree{\ensuremath{\twofun{F^{JP}}}}  %free functor from FPFib to JPComp
\def\JTFree{\ensuremath{\twofun{F^T}}}  %free functor from JComp to JTComp

\def\C#1#2{\ensuremath{c^{#1}_{#2}}}  %frecce c_i^k
\def\Kdom#1{\ensuremath{\dom\cdot#1}}   %dominio di \chi
\newcommand{\jcomp}[1]{\chi^{#1}} %comprensione di Lawvere
\def\JFC#1{\ensuremath{\jcomp{\JFib{#1}}}}  %funtore comprensione libero su fibrazione
\newcommand{\weak}[2]{\mathrm{w}_{#1}^*#2}  %weakening di 2 lungo comprensione di 1
\newcommand{\weakar}[2]{\mathrm{w}_{#1}#2}  %weakening arrow di 2 lungo comprensione di 1
\newcommand{\gel}[1]{\mathrm{g}_{#1}}  %elemento generico su 1
\newcommand{\var}[1]{x}  %variabile generica

\newcommand{\Tr}{\OPR{Tr}}  %category pre-quotient
\newcommand{\Trfib}[1]{\OPR{Trfib} #1}  %category pre-quotient
\newcommand{\Const}{\OPR{Const}}  %constants functor
\newcommand{\ConstCat}{\OPR{ConstC}}  %constants functor
\newcommand{\ConstFib}[1]{\overline{#1}}  %functor from constcat to base
\newcommand{\triang}{tr}  %constant to triangles functor
\newcommand{\incl}{incl}  %constant to old morphisms
\newcommand{\LEFree}{\twofun{F^{LE}}}  
\newcommand{\LEFib}[1]{{#1}^{\twofun{LE}}}
\newcommand{\LEFun}[1]{{#1}^{\twofun{LE}}}
\newcommand{\LEDom}[2]{{#1}^{\twofun{LE}}_{#2}}
\newcommand{\LEDomE}[1]{\LEDom{\ct{E}}{#1}}
\newcommand{\LEun}[1]{\twont{\eta}^{\twofun{LE}}_{#1}}
\newcommand{\LEcoun}[1]{\twont{\epsilon}^{\twofun{LE}}_{#1}}
\newcommand{\quot}[1]{\OPR{Q}^{#1}}  %quotient towards LEFib
\newcommand{\semtocat}{\OPR{F}}  %free functor from semicategories to categories ??

%% Theorems
\newtheorem*{notation}{Notation}


\newcommand{\fibre}[3][]{#2\ifblank{#1}{}{^{#1}}_{#3}} 
